\documentclass{amsart}
% For dealing with references we use the comment environment
\usepackage{verbatim}
\newenvironment{reference}{\comment}{\endcomment}
%\newenvironment{reference}{}{}
\newenvironment{slogan}{\comment}{\endcomment}
\newenvironment{history}{\comment}{\endcomment}

% For commutative diagrams we use Xy-pic
\usepackage[all]{xy}

% We use 2cell for 2-commutative diagrams.
\xyoption{2cell}
\UseAllTwocells

% We use multicol for the list of chapters between chapters
\usepackage{multicol}

% This is generally recommended for better output
\usepackage{lmodern}
\usepackage[T1]{fontenc}

% For cross-file-references

% Package for hypertext links:
\usepackage{hyperref}

% For any local file, say "hello.tex" you want to link to please

% Theorem environments.
%
\theoremstyle{plain}
\newtheorem{theorem}[subsection]{Theorem}
\newtheorem{proposition}[subsection]{Proposition}
\newtheorem{lemma}[subsection]{Lemma}

\theoremstyle{definition}
\newtheorem{definition}[subsection]{Definition}
\newtheorem{example}[subsection]{Example}
\newtheorem{exercise}[subsection]{Exercise}
\newtheorem{situation}[subsection]{Situation}

\theoremstyle{remark}
\newtheorem{remark}[subsection]{Remark}
\newtheorem{remarks}[subsection]{Remarks}

\numberwithin{equation}{subsection}

% Macros
%
\def\lim{\mathop{\mathrm{lim}}\nolimits}
\def\colim{\mathop{\mathrm{colim}}\nolimits}
\def\Spec{\mathop{\mathrm{Spec}}}
\def\Hom{\mathop{\mathrm{Hom}}\nolimits}
\def\Ext{\mathop{\mathrm{Ext}}\nolimits}
\def\SheafHom{\mathop{\mathcal{H}\!\mathit{om}}\nolimits}
\def\SheafExt{\mathop{\mathcal{E}\!\mathit{xt}}\nolimits}
\def\Sch{\mathit{Sch}}
\def\Mor{\mathop{\mathrm{Mor}}\nolimits}
\def\Ob{\mathop{\mathrm{Ob}}\nolimits}
\def\Sh{\mathop{\mathit{Sh}}\nolimits}
\def\NL{\mathop{N\!L}\nolimits}
\def\CH{\mathop{\mathrm{CH}}\nolimits}
\def\proetale{{pro\text{-}\acute{e}tale}}
\def\etale{{\acute{e}tale}}
\def\QCoh{\mathit{QCoh}}
\def\Ker{\mathop{\mathrm{Ker}}}
\def\Im{\mathop{\mathrm{Im}}}
\def\Coker{\mathop{\mathrm{Coker}}}
\def\Coim{\mathop{\mathrm{Coim}}}

% Boxtimes
%
\DeclareMathSymbol{\boxtimes}{\mathbin}{AMSa}{"02}

%
% Macros for moduli stacks/spaces
%
\def\QCohstack{\mathcal{QC}\!\mathit{oh}}
\def\Cohstack{\mathcal{C}\!\mathit{oh}}
\def\Spacesstack{\mathcal{S}\!\mathit{paces}}
\def\Quotfunctor{\mathrm{Quot}}
\def\Hilbfunctor{\mathrm{Hilb}}
\def\Curvesstack{\mathcal{C}\!\mathit{urves}}
\def\Polarizedstack{\mathcal{P}\!\mathit{olarized}}
\def\Complexesstack{\mathcal{C}\!\mathit{omplexes}}
% \Pic is the operator that assigns to X its picard group, usage \Pic(X)
% \Picardstack_{X/B} denotes the Picard stack of X over B
% \Picardfunctor_{X/B} denotes the Picard functor of X over B
\def\Pic{\mathop{\mathrm{Pic}}\nolimits}
\def\Picardstack{\mathcal{P}\!\mathit{ic}}
\def\Picardfunctor{\mathrm{Pic}}
\def\Deformationcategory{\mathcal{D}\!\mathit{ef}}

\setlength{\emergencystretch}{3em}
\begin{document}
\title[Stacks Project, Tag 09NX]{584 --- Every reduced separated Noetherian scheme of dimension one admits an ample line bundle}
\maketitle
\noindent Copyright (C) 2005--2025 Johan de Jong. Stacks Project authors. Source: \href{https://stacks.math.columbia.edu/tag/09NX}{Tag 09NX}.

\noindent Editorial difficulty note (not author text). Difficulty H2: Ample bundles on the irreducible components must be glued compatibly near all component intersections. The proof constructs a common affine section neighbourhood of those intersections, then uses conductor ideals and products of component sections to produce affine principal neighbourhoods at every remaining point..

\noindent Editorial provenance note (not author text). History: excerpt from revision \texttt{a0444\allowbreak 6e57e\allowbreak c1fbc\allowbreak 252a8\allowbreak 71afc\allowbreak ec775\allowbreak 2fb28\allowbreak 07b14}, varieties.tex, lines 8079--8135. Reference presentation was adapted; original-author context and core support blocks were retained or added. The mathematical wording is unchanged. Licensed under GNU FDL 1.2 or later; no invariant sections or cover texts. See ../licenses/COPYING. Context blocks reproduce author definitions and supporting results; they are not additional counted problems.

\begin{definition}
\label{properties-definition-ample}
\begin{reference}
\href{https://stacks.math.columbia.edu/bibliography}{Bibliography: EGA (II Definition 4.5.3)}
\end{reference}
Let $X$ be a scheme.
Let $\mathcal{L}$ be an invertible $\mathcal{O}_X$-module.
We say $\mathcal{L}$ is {\it ample} if
\begin{enumerate}
\item $X$ is quasi-compact, and
\item for every $x \in X$ there exists an $n \geq 1$
and $s \in \Gamma(X, \mathcal{L}^{\otimes n})$ such
that $x \in X_s$ and $X_s$ is affine.
\end{enumerate}
\end{definition}

\begin{lemma}
\label{lemma-dim-1-noetherian-reduced-separated-has-ample}
Let $X$ be a Noetherian reduced separated scheme of dimension $1$.
Then $X$ has an ample invertible sheaf.
\end{lemma}

\begin{proof}
Let $Z_i$, $i = 1, \ldots, n$ be the irreducible components of $X$.
We view these as reduced closed subschemes of $X$.
By Lemma \href{https://stacks.math.columbia.edu/tag/09ND}{lemma dim 1 noetherian integral separated has ample (Tag 09ND)}
there exist ample invertible sheaves $\mathcal{L}_i$ on $Z_i$.
Set $T = \bigcup_{i \not = j} Z_i \cap Z_j$. As $X$ is Noetherian
of dimension $1$, the set $T$ is finite and consists of closed
points of $X$. For each $i$ we may, possibly after replacing
$\mathcal{L}_i$ by a power, choose $s_i \in \Gamma(Z_i, \mathcal{L}_i)$
such that $(Z_i)_{s_i}$ is affine and contains $T \cap Z_i$, see
Properties, Lemma \href{https://stacks.math.columbia.edu/tag/09NV}{lemma ample finite set in principal affine (Tag 09NV)}.

\medskip\noindent
By Lemma \href{https://stacks.math.columbia.edu/tag/09NE}{lemma glue invertible sheaves (Tag 09NE)} we can find an invertible sheaf
$\mathcal{L}$ on $X$ and $s \in \Gamma(X, \mathcal{L})$
such that $(\mathcal{L}, s)|_{Z_i} = (\mathcal{L}_i, s_i)$.
Observe that $X_s$ contains $T$ and is set theoretically equal
to the affine closed subschemes $(Z_i)_{s_i}$. Thus it is affine by
Limits, Lemma \href{https://stacks.math.columbia.edu/tag/09NL}{lemma affines glued in closed affine (Tag 09NL)}.
To finish the proof, it suffices to find for every $x \in X$, $x \not \in T$
an integer $m > 0$ and a section $t \in \Gamma(X, \mathcal{L}^{\otimes m})$
such that $X_t$ is affine and $x \in X_t$. Since $x \not \in T$
we see that $x \in Z_i$ for some unique $i$, say $i = 1$.
Let $Z \subset X$ be the reduced closed subscheme whose underlying
topological space is $Z_2 \cup \ldots \cup Z_n$.
Let $\mathcal{I} \subset \mathcal{O}_X$ be the ideal
sheaf of $Z$. Denote that $\mathcal{I}_1 \subset \mathcal{O}_{Z_1}$
the inverse image of this ideal sheaf under the inclusion
morphism $Z_1 \to X$. Observe that
$$
\Gamma(X, \mathcal{I}\mathcal{L}^{\otimes m}) =
\Gamma(Z_1, \mathcal{I}_1 \mathcal{L}_1^{\otimes m})
$$
see Remark \href{https://stacks.math.columbia.edu/tag/09NW}{remark conductor (Tag 09NW)}. Thus it suffices to find $m > 0$
and $t \in \Gamma(Z_1, \mathcal{I}_1 \mathcal{L}_1^{\otimes m})$
with $x \in (Z_1)_t$ affine. Since $\mathcal{L}_1$ is ample
and since $x$ is not in $Z_1 \cap T = V(\mathcal{I}_1)$
we can find a section
$t_1 \in \Gamma(Z_1, \mathcal{I}_1 \mathcal{L}_1^{\otimes m_1})$
with $x \in (Z_1)_{t_1}$, see Properties, Proposition
\href{https://stacks.math.columbia.edu/tag/01Q3}{proposition characterize ample (Tag 01Q3)}.
Since $\mathcal{L}_1$ is ample we can find a section
$t_2 \in \Gamma(Z_1, \mathcal{L}_1^{\otimes m_2})$
with $x \in (Z_1)_{t_2}$ and $(Z_1)_{t_2}$ affine, see
Properties, Definition \ref{properties-definition-ample}.
Set $m = m_1 + m_2$ and $t = t_1 t_2$. Then
$t \in \Gamma(Z_1, \mathcal{I}_1 \mathcal{L}_1^{\otimes m})$
with $x \in (Z_1)_t$ by construction and
$(Z_1)_t$ is affine by Properties, Lemma
\href{https://stacks.math.columbia.edu/tag/01PV}{lemma affine cap s open (Tag 01PV)}.
\end{proof}
\end{document}
